\section{Theorem 1: Capacity-Induced Module Curvature Disparity}
\label{sec:thm1}

This section establishes our first main result: the module curvature
disparity $\mathcal{D}_{\mathrm{curv}}$
(Definition~\ref{def:dcurv}) grows with the spatial module's width.
The result is the geometric foundation for the dynamical analysis of
Section~\ref{sec:thm2}: it identifies a capacity-dependent mechanism
capable of generating disparate module learning rates. The dynamical
starvation claim itself is stated separately in
Theorem~\ref{thm:twoscale}, under its own explicit hypotheses.

\subsection{Main result}

\begin{theorem}[Capacity-Induced Module Curvature Disparity]
\label{thm:hessian}
Consider the compositional model $F_{\thetap,\phip} = g_\phip \circ
f_\thetap$ with fixed-dimensional linear spectral module $f_\thetap$
(Assumption~\ref{ass:linear}) and a family of spatial networks
$g_\phip^{(M)}$ indexed by characteristic width $M$
(channels for a CNN, embedding dimension for a ViT, hidden width for
an MLP). Under
Assumptions~\ref{ass:linear}, \ref{ass:subg},
and~\ref{ass:inputlip}, together with the head
hypothesis of Lemma~\ref{lem:phi_scaling} below (a dense classifier
head over penultimate features of dimension $\Theta(M)$ that are
correlated across inputs at initialization), there exist
constants $c_1, c_2 > 0$ depending on the data distribution, the
initialization scheme, the bottleneck dimension $K$, the input
spectral dimension $S$, and the architecture class of $g_\phip$
(depth, activation, head dimensions)---but not on $M$ or on the
spatial parameter count $\Cg$---such that
\[
    \mathcal{D}_{\mathrm{curv}}^{(M)}
    \;=\; \frac{\lambda_{\max}(G^{(M)}_{\phip\phip})}
                {\lambda_{\max}(G_{\thetap\thetap})}
    \;\geq\; c_1 \cdot M - c_2
\]
in expectation over random Kaiming (variance-scaled Gaussian)
initialization. Equivalently, for architecture families with
$\Cg = \Theta(M^2)$,
\[
    \mathcal{D}_{\mathrm{curv}}^{(M)}
    \;\geq\; c_1' \cdot \sqrt{\Cg / \Cf} - c_2.
\]
\end{theorem}

\begin{remark}[Interpretation]
Theorem~\ref{thm:hessian} states that as the spatial network widens,
its top Gauss-Newton eigenvalue grows without bound while the
spectral module's top Gauss-Newton eigenvalue remains capped by a
$\Cg$-independent constant. Curvature increasingly concentrates in
the downstream module. This is a purely geometric fact about the
Hessian; the dynamical consequence---that joint training pins
$\thetap$ near its initialization---is stated separately in
Theorem~\ref{thm:twoscale}, under its own hypotheses. The disparity
measures the downstream module's \emph{readiness to move}---the
loss curvature its parameters command at initialization---not the
usefulness of its features: it is present before any feature has
been learned, and the width-linear spike can be exhibited even for
inputs carrying no information
(Remark~\ref{rem:instance_interpretation}, supplement). That the
geometry favours the spatial module \emph{independently} of feature
quality is precisely what makes it a shortcut hazard.
\end{remark}

The proof decomposes the disparity bound into separate bounds on the
top eigenvalues of the two diagonal blocks. We state the supporting
lemmas before presenting the proof.

\subsection{Supporting lemmas}

\begin{lemma}[Spatial block scaling]
\label{lem:phi_scaling}
Let $M$ denote the characteristic width of $g_\phip$ (channel count
for a CNN, embedding dimension for a ViT, hidden layer width for an
MLP), let $N_{\mathrm{cls}} \geq 2$, and suppose $g_\phip$ ends in a dense
classifier head over penultimate features
$h_n \in \mathbb{R}^{M_h}$ with $M_h = \Theta(M)$ whose
softmax-weighted mean Gram
$S_h^{p} := N^{-1}\sum_n p^{(n)}_{\min} h_n h_n^\top$
(Equation~\eqref{eq:feature_gram}) satisfies
\[
    \lambda_{\max}(S_h^{p}) \;\geq\; q_h M_h
    \quad\text{with probability } \geq 1 - \delta_h
\]
over the random initialization---where $p^{(n)}_{\min}$ denotes the
smallest softmax probability at pixel $n$---for constants $q_h > 0$
and $\delta_h < 1$ independent of $M$ (the \emph{head hypothesis}:
the features carry a common direction of $\Theta(M)$ energy across
inputs, on pixels whose initial predictions are non-degenerate. It
is a cross-input \emph{correlation} condition, not a per-vector norm
condition, and is closely related to the quantity $\kappa_2$
carrying the width-linear term in Karakida et al.'s
law---see the supplement, where a fully proved instance is also
given). Then the spatial Gauss-Newton
block of~\eqref{eq:ggn} satisfies
\[
    \lambda_{\max}(G_{\phip\phip})
    \;=\; \Omega(M)
    \qquad \parbox{15em}{\centering in expectation over random Kaiming
    (variance-scaled Gaussian) initialization.}
\]
The logit Gram $J_\phip^\top J_\phip \succeq G_{\phip\phip}$
inherits the bound unconditionally; the residual Gram
$J_\phip^{\mathrm{res}\top} J_\phip^{\mathrm{res}}$ satisfies it
under the analogous hypothesis on the $p^2$-weighted Gram
$S_h^{p^2}$ (supplement). For architecture families with
$\Cg = \Theta(M^2)$ (fixed depth with at least two consecutive
width-proportional layers), equivalently
$\lambda_{\max}(G_{\phip\phip}) = \Omega(\sqrt{\Cg})$.
\end{lemma}

\begin{remark}
Lemma~\ref{lem:phi_scaling} is proved in the supplement
(Section~\ref{sec:proof_thm1}) by an explicit witness: the rank-one
head-weight perturbation $V = c\, u^\top$, where $c$ is a unit class
contrast and $u$ the dominant direction of the feature ensemble.
Evaluating the curvature along $V$ gives the deterministic
inequality $\lambda_{\max}(G_{\phip\phip}) \geq
\lambda_{\max}(S_h^{p})$, and the head hypothesis supplies
$\lambda_{\max}(S_h^{p}) = \Omega(M)$. The
width-linear law agrees with the mean-field analysis of Karakida et
al.~\cite{karakida2019universal}: their Theorem~4 gives
$\lambda_{\max} = \Theta(M)$ with sharp constant for
proportional-width fully connected networks under squared loss, and
their Theorem~1 shows the mean eigenvalue simultaneously shrinks as
$\mathcal{O}(1/M)$---curvature concentrates in a vanishing fraction
of directions as networks widen. The scaling variable is width, not
parameter count: for architecture families with $\Cg = \Theta(M^2)$,
the equivalent parameter-count statement is
$\lambda_{\max} = \Omega(\sqrt{\Cg})$. The head hypothesis covers
per-pixel dense heads on CNN and ViT modules as well; the sharpness
of the width-linear rate is tested empirically in
Section~\ref{sec:experiments}.
\end{remark}

\begin{lemma}[Spectral block operator-norm cap]
\label{lem:opcap}
Under Assumptions~\ref{ass:linear} (linear per-pixel $f_\thetap$),
\ref{ass:subg} (bounded input second moment), and~\ref{ass:inputlip}
(bounded input Jacobian
$\|\partial g_\phip / \partial Z\|_{\mathrm{op}} \leq \widetilde{L}$
uniformly in $\Cg$), the spectral logit Jacobian and Gauss-Newton
block satisfy
\[
    \lambda_{\max}(G_{\thetap\thetap})
    \;\leq\; \|J_\thetap\|_{\mathrm{op}}^{2}
    \;\leq\; \widetilde{L}^{2} \cdot \lambda_{\max}(\Sigma_X)
    \;=:\; C_\thetap,
\]
where $C_\thetap$ depends only on the classifier's input Lipschitz
constant $\widetilde{L}$ and the top eigenvalue of the mean per-pixel
input Gram $\Sigma_X$ of~\eqref{eq:sigma_x}---not on the spatial
width $M$, the spatial parameter count $\Cg$, or the dataset size
$N$. The same bound holds for the residual-Gram convention, since
$H_\tau^2 \preceq H_\tau \preceq I$.
\end{lemma}

\begin{remark}
Lemma~\ref{lem:opcap} follows from the chain-rule identity
$J_\thetap = J_Z \cdot (\partial Z / \partial \thetap)$ for the logit
Jacobian, combined with $\|J_Z\|_{\mathrm{op}} \leq \widetilde{L}$
(Assumption~\ref{ass:inputlip}) and the exact form of $\partial Z /
\partial \thetap$ under linearity of $f_\thetap$
(Assumption~\ref{ass:linear}, giving the per-pixel factor $I_K
\otimes x_n^\top$, whose $N$-normalized stacking has
$\|\partial Z / \partial \thetap\|_{\mathrm{op}}^{2} =
\lambda_{\max}(\Sigma_X)$ exactly, by
$I_K \otimes \Sigma_X$). Passing from the logit Jacobian to the
Gauss-Newton block or the residual Jacobian only inserts factors of
$H_\tau \preceq I$ and can only decrease the bound. The bound is on
the operator norm (equivalently, the top eigenvalue), which is
well-behaved under rank-deficiency of the block---avoiding the
smallest-eigenvalue and Kronecker-factorisation machinery discussed
in Section~\ref{sec:proof_thm1}.
\end{remark}

\subsection{Proof of Theorem~\ref{thm:hessian}}

\begin{proof}[Proof of Theorem~\ref{thm:hessian}]
By Lemma~\ref{lem:phi_scaling}, in expectation over random Kaiming
(variance-scaled Gaussian) initialization,
\[
    \lambda_{\max}(G^{(M)}_{\phip\phip}) \;\geq\; c_\phi \cdot M
\]
for a constant $c_\phi > 0$ independent of $M$ (depending on the
initialization scheme, the activation, the head dimensions, and the
data distribution; supplement Section~\ref{sec:proof_thm1}). By
Lemma~\ref{lem:opcap},
\[
    \lambda_{\max}(G_{\thetap\thetap})
    \;\leq\; \|J_\thetap\|^2_{\mathrm{op}} \;\leq\; C_\thetap
\]
for a constant $C_\thetap > 0$ independent of $M$. The cap is
deterministic---it holds for every realization of the
initialization---so it may be divided through inside the
expectation:
\[
    \E\!\left[\mathcal{D}_{\mathrm{curv}}^{(M)}\right]
    \;=\; \E\!\left[\frac{\lambda_{\max}(G^{(M)}_{\phip\phip})}
                {\lambda_{\max}(G_{\thetap\thetap})}\right]
    \;\geq\; \frac{\E\!\left[\lambda_{\max}(G^{(M)}_{\phip\phip})\right]}
                  {C_\thetap}
    \;\geq\; \frac{c_\phi \cdot M}{C_\thetap}
    \;=\; c_1 \cdot M
\]
with $c_1 := c_\phi / C_\thetap$. The witness route delivers the
bound with $c_2 = 0$; the subtracted constant is retained in the
statement so that sharper asymptotic constants, valid only for
large $M$, may be substituted. Under
$\Cg = \Theta(M^2)$ (families with at least two consecutive
width-proportional layers), $M = \Theta(\sqrt{\Cg})$
and the ratio form $c_1' \cdot \sqrt{\Cg / \Cf}$ follows on multiplying
and dividing by $\sqrt{\Cf}$. This completes the proof.
\end{proof}

\subsection{Implications}

Theorem~\ref{thm:hessian} establishes that curvature concentrates in
the downstream spatial module as its width grows. We highlight two
implications, and separate them explicitly from the dynamical
consequences of Section~\ref{sec:thm2}.

\paragraph{Geometric mechanism, not a dynamical claim.}
$\mathcal{D}_{\mathrm{curv}} \gg 1$ identifies a
\emph{potential} for disparate module learning rates: the spatial
Gauss-Newton block admits directions of arbitrarily large curvature
while the spectral block does not. Whether joint gradient descent
actually produces disparate effective learning rates depends on an
additional dynamical hypothesis (that the shortcut pathway fits the
residual at rate $\mu$; Assumption~\ref{ass:residual}), which
Section~\ref{sec:thm2} introduces and Experiment~1.2 tests.
Theorem~\ref{thm:hessian} does not,
by itself, prove that $\thetap$ is pinned to initialization.
In particular, $\mathcal{D}_{\mathrm{curv}}$ concerns the \emph{top}
of the spatial block's spectrum, while the fitting rate $\mu$ of
Section~\ref{sec:thm2} is governed by how fast the residual can be
driven down---a property of other parts of the spatial spectrum. The
two are linked only under spectral-shape assumptions we do not make;
the theorems are deliberately independent, and
Experiments~1.1--1.2 test them separately.

\paragraph{The disparity is invariant to shared step-size tuning.}
The disparity bound depends only on architectural parameters, not on
the choice of optimizer or the shared step-size schedule---so no
choice of shared learning rate changes the \emph{geometry}; whether
any shared-step-size method escapes the \emph{dynamical} consequence
is a separate question we treat as a conjecture below. Adaptive
methods (Adam, RMSProp) rescale per-coordinate gradient magnitudes
but do not alter the underlying eigenvalue disparity between the two
blocks. Two-timescale optimizers with explicitly different learning
rates per module (Heusel et al.~\cite{heusel2017gans}; Borkar
2008~\cite{borkar2008stochastic} Ch.~6) are the natural remedy at the
optimizer level. We discuss this in Section~\ref{sec:discussion}.

\subsection{Relation to classical ill-conditioning}
\label{sec:classical_conditioning}

Ill-conditioned optimization has classical causes with classical
remedies, and it is natural to ask whether they apply here. They do
not, because the disparity of Theorem~\ref{thm:hessian} is driven by
the \emph{architecture}, not by the data or the parametrization of a
single block.

\paragraph{Feature normalization.}
Scale asymmetry or collinearity among input features degrades the
conditioning of $\Sigma_X$ and is classically cured by
standardizing (e.g.\ mean-centering) the inputs. Here the cure has
a sharper, quantitative role: Lemma~\ref{lem:opcap} places
$\lambda_{\max}(\Sigma_X)$ in the spectral cap, so input
transformations that deflate the top of the input spectrum
\emph{raise} the realized disparity at every width---normalization
moves the constants. What no input transformation can touch is the
width-scaling of the spatial block (Lemma~\ref{lem:phi_scaling}):
capacity moves the rate. On raw nonnegative spectra, whose second
moment is dominated by the shared mean shape, the cap is large and
the realized disparity correspondingly small; mean-centering is
predicted to restore it (Section~\ref{sec:real}).

\paragraph{Weight decay.}
Adding $\alpha \norm{w}^2$ to the loss shifts the Hessian by
$2\alpha I$, lifting every eigenvalue by $2\alpha$ and famously
repairing classical condition numbers. But the shift acts on
\emph{both} blocks equally: the disparity becomes
$(\lambda_{\max}(G_{\phip\phip}) + 2\alpha) /
(\lambda_{\max}(G_{\thetap\thetap}) + 2\alpha)$, which remains
$\Theta(M)$ for any fixed $\alpha$ as $M$ grows. Repairing the
disparity would require $2\alpha \sim \lambda_{\max}(G_{\phip\phip})$,
a regularization strength that would dominate the data term and
prevent fitting. Consistent with this, Experiment~1.6 finds that
Pezeshki's Spectral Decoupling---an $L_2$ penalty on logits designed
to combat saturation---does not reduce the joint-training collapse,
while freezing does.

\paragraph{Adaptive optimizers.}
Adam and RMSProp precondition per-coordinate gradient magnitudes with
a diagonal rescaling. The raw block spectra are
optimizer-invariant, but the \emph{dynamics} under a diagonal
preconditioner $P$ are governed by the preconditioned spectra of
$P^{-1/2} G P^{-1/2}$ and need not respect the disparity: under
Adam, per-coordinate steps are near-constant in gradient magnitude,
so small spectral gradients do not imply small spectral movement.
The gradient-flow statements of Section~\ref{sec:thm2} are
accordingly scoped to (stochastic) gradient descent; joint-vs-frozen
gaps observed under Adam are evidence about the phenomenon, not
applications of the theorem. Empirically, the
canonical dynamics experiment (Experiment~1.2) is run under Adam and
the joint-vs-frozen gap persists at every width; we treat the stronger
claim---that no shared-learning-rate adaptive method can remove the
pathology---as a conjecture rather than a theorem. Optimizers with
explicitly different learning rates per module (two-timescale update
rules; Heusel et al.~\cite{heusel2017gans}) are the natural remedy at
the optimizer level and correspond, in the limit, to our freezing
prescription.

\begin{remark}[Parameterization scope]
\label{rem:param_scope}
$\mathcal{D}_{\mathrm{curv}}$ is not invariant to per-module
reparameterization, and Theorem~\ref{thm:hessian} is a statement
about the standard variance-scaled (Kaiming) parameterization---the
one used in the pipelines we analyze and in all our experiments.
Under alternative parameterizations that rescale layers with width,
such as $\mu$P~\cite{yang2021tensor}, both the initialization
magnitudes and the effective per-layer learning rates change, and
the width scaling of the blocks would need to be re-derived; we do
not claim the disparity law transfers. Our two uses of $\mu$P
elsewhere are limited to a different role: as one route to keeping
the input-Jacobian bound (Assumption~\ref{ass:inputlip}) uniform in
width during training. Whether principled reparameterization can
neutralize the disparity---rather than merely re-expressing it---is
an open question adjacent to the per-module learning-rate remedies
discussed in Section~\ref{sec:discussion}.
\end{remark}

\begin{remark}[Small eigenvalues: irrelevance vs.\ starvation]
\label{rem:irrelevance_vs_starvation}
A parameter direction can carry vanishing curvature for two very
different reasons: (a) the direction genuinely does not affect
predictions on the task at hand (``irrelevance''), in which case zero
gradient is the correct outcome and the parameter should not be
learned; or (b) the direction \emph{would} affect predictions if the
optimizer could reach it, but the compositional geometry buries the
gradient signal behind the bottleneck (``starvation''). The
Gauss-Newton bounds of Theorem~\ref{thm:hessian} alone cannot
distinguish between these interpretations: both produce
$\Cg$-independent small curvature in the spectral block.

The synthetic experiments in Section~\ref{sec:experiments} are
designed to distinguish these hypotheses by construction. In
Experiment~1.3 we generate data where each of the two modalities alone
carries information sufficient to solve the task (single-modality
baselines reach $\approx 0.53$ on a binary classification, above the
$0.50$ chance floor). Joint training on the same data does not merely
under-use the spectral channel: it undershoots the weaker of the two
single-modality baselines. The gap is modest in absolute terms on this
synthetic problem---an important honest caveat---but its direction is
inconsistent with interpretation~(a), under which joint training
would at worst match the stronger single-modality baseline.

For real spectroscopic data, distinguishing (a) from (b) requires
demonstrating that a frozen pretrained spectral encoder outperforms
one that has been jointly trained end-to-end on the same task. We
report such experiments on FTIR and QCL tissue classification in
Section~\ref{sec:real}. Our present study compares only two
concrete mitigation strategies (freezing the spectral module, and
Pezeshki's Spectral
Decoupling~\cite{pezeshki2021gradient}); broader comparisons against
modality-specific learning rates, gradient surgery, GradNorm-style
loss balancing, and related interventions are important future work.
\end{remark}

\subsection{Empirical support}

Experiment~1.1 in Section~\ref{sec:experiments} measures
$\mathcal{D}_{\mathrm{curv}}^{(M)}$ directly across a range of
spatial widths and confirms monotonic growth with $M$: the spatial
top eigenvalue grows $74\times$ over the sweep while the spectral
top eigenvalue stays flat, driving the disparity into the thousands.
The measured log--log slope of $\lambda_\phi$ against $\Cg$ is
$0.70 \pm 0.05$, below the slope-$1.0$ prediction that applies to
the tested architecture family (whose $\Cg = \Theta(M)$); the
deficit and its candidate causes are reported in
Section~\ref{sec:experiments}. The qualitative
prediction---unbounded spatial-block growth against a capped
spectral block---is unambiguous in the data.

\begin{remark}[Defense against linear triviality]
\label{rem:not_pca}
A natural objection is that, since $f_\thetap$ is linear, the analysis
must reduce to principal component analysis (PCA) of the input. This
is not the case. PCA finds directions of maximum input variance
\emph{without reference to labels}. At initialization, however, the
distinction is limited, and we state this candidly: the
Gauss--Newton blocks carry no label information there (the softmax
weights are near-uniform), and $G_{\thetap\thetap}$ is, by
Lemma~\ref{lem:opcap}, an input-second-moment object propagated
through a label-free random network. The defense against
PCA-triviality is therefore a statement about \emph{training}: along
the trajectory the supervised residual shapes
$\partial \hat{y}/\partial Z$ and hence the curvature, and the
pathology we identify is precisely that this supervised signal
fails to reach $\thetap$---not that $f_\thetap$ is incapable of
representing useful features.
\end{remark}
